\subsection{代数的逆向极限与正向极限}

设 I 是一个预序集， \((\mathbf{A}_{i},\phi_{ij})\) 是以 I 为指标集的交换环的逆向系统。设 \((\mathbf{E}_{i},f_{ij})\) 是以 I 为指标集的 \(A_{i}\) -模的逆向系统（II，§ 6，第 1 号），并进一步假设每个 \(E_{i}\) 具有 \(A_{i}\) -代数结构，且对 \(i\leqslant j,f_{ij}\) 是代数的 \(A_{j}\) -同态（相对于 \(\phi_{ij}\) ）（第 5 号）。设 \(A=\varprojlim A_{i}\) 与 \(E=\varprojlim E_{i}\) ，后者具有 A-模结构，即诸 \(\mathbf{A}_i\) -模 \(\mathbf{E}_i\) 的结构的逆向极限（II，§ 6，第 1 号）；立即可验证， \(\mathbf{E}\) 上的合成律——把诸 \(\mathbf{E}_i\) 视为乘法下的广群（I，§ 10，第 1 号）而取其逆向极限——连同 \(\mathbf{E}\) 上的 A-模结构，在 \(\mathbf{E}\) 上定义了一个 A-代数结构； \((\mathbf{E}_i, f_{ij})\) 称为 \(\mathbf{A}_i\) -代数的逆向系统，而 A-代数 \(\mathbf{E}\) 称为其逆向极限。若 \(f_i: \mathbf{E} \to \mathbf{E}_i\) , \(\phi_i: \mathbf{A} \to \mathbf{A}_i\) 为典范映射，则 \(f_i\) 是代数的 A-同态（相对于 \(\phi_i\) ）。若诸 \(\mathbf{E}_i\) 结合（相应地，交换），则 \(\mathbf{E}\) 亦然；若每个 \(\mathbf{E}_i\) 具有单位元 \(e_i\) 且 \(f_{ij}(e_j) = e_i\) （对 \(i \leqslant j\) ），则 \(e = (e_i)\) 是代数 \(\mathbf{E}\) 的单位元。

设 \((\mathbf{E}_{i}^{\prime}, f_{ij}^{\prime})\) 是 \(A_{i}\) -代数的另一个逆向系统，且对所有 i 令 \(u_{i}: E_{i} \to E_{i}^{\prime}\) 为 \(A_{i}\) -代数同态，这些映射构成一个逆向系统；则 \(u = \lim u_{i}\) 是一个 A-代数同态。

现在假设所有的 \(A_{i}\) 都等于同一个交换环 A，且诸 \(\phi_{ij}\) 等于 \(Id_{A}\) ，于是 \(E = \lim E_{i}\) 是一个 A-代数。设 F 是一个 A-代数，且对所有 \(i \in I\) ，设 \(u_{i}: F \to E_{i}\) 是一个 A-代数同态，使得 \((u_{i})\) 是一个映射的逆向系统；则 \(u = \lim u_{i}\) 是代数 F 到代数 E 中的同态。反之，对每个 A-代数同态 \(v: F \to E\) ，族 \(v_{i} = f_{i} \circ v\) 是一个 A-代数同态的逆向系统，使得 \(v = \lim v_{i}\) 。此外，由于记 \(\bar{f}_{ij} = \text{Hom}(1_{\text{F}}, f_{ij})\) , \((\text{Hom}_{\text{A-alg.}}(\text{F}, \text{E}_i), \bar{f}_{ij})\) 显然是一个集合的逆向系统，可见上述评注也可以表述为：典范映射 \(v \mapsto (f_i \circ v)\) 是一个双射

\[
l _ {\mathrm{F}}: \operatorname{Hom} _ {\mathrm{A-alg.}} (\mathrm{F}, \varprojlim \mathrm{E} _ {i}) \rightarrow \varprojlim \operatorname{Hom} _ {\mathrm{A-alg.}} (\mathrm{F}, \mathrm{E} _ {i}).
\]

此外，对于每一个 A-代数同态 \(w: \mathbf{F} \to \mathbf{F}'\) ，

\[
\bar {w} _ {i} = \operatorname{Hom} (w, 1 _ {\mathrm{E} _ {i}}): \operatorname{Hom} _ {\mathrm{A-alg.}} (\mathrm{F} ^ {\prime}, \mathrm{E} _ {i}) \rightarrow \operatorname{Hom} _ {\mathrm{A-alg.}} (\mathrm{F}, \mathrm{E} _ {i})
\]

构成映射的逆向系统，且图表

\[
\begin{array}{c} \operatorname{Hom} _ {\mathrm{A-alg.}} (\mathrm{F} ^ {\prime}, \lim \mathrm{E} _ {i}) \xrightarrow {l _ {\mathbf {F} ^ {\prime}}} \lim \operatorname{Hom} _ {\mathrm{A-alg.}} (\mathrm{F} ^ {\prime}, \mathrm{E} _ {i}) \\ \operatorname{Hom} (w, 1 _ {\mathbf {E}}) \Biggl \downarrow \qquad \qquad \qquad \qquad \qquad \qquad \qquad \qquad \qquad \qquad \qquad \qquad \qquad \qquad \qquad \qquad \qquad \qquad \qquad \qquad \qquad \qquad \qquad \qquad \qquad \qquad \qquad \qquad \qquad \qquad \qquad \qquad \qquad \qquad \\ \operatorname{Hom} _ {\mathrm{A-alg.}} (\mathrm{F}, \lim \mathrm{E} _ {i}) \xrightarrow {l _ {\mathbf {F}}} \lim \operatorname{Hom} _ {\mathrm{A-alg.}} (\mathrm{F}, \mathrm{E} _ {i}) \end{array}
\]

是交换的。

现在假设 I 是右有向的。考虑交换环的正向系统 \((\mathbf{A}_i, \phi_{ji})\) 与正向系统 \((\mathbf{E}_i, f_{ji})\) （ \(\mathbf{A}_i\) -模的），均以 I 为指标集；假设每个 \(\mathbf{E}_i\) 具有 \(\mathbf{A}_i\) -代数结构，且对 \(i \leqslant j\) ， \(f_{ji}\) 是代数的 \(\mathbf{A}_i\) -同态（相对于 \(\phi_{ji}\) ）（第 5 号）。设 \(\mathbf{A} = \varinjlim \mathbf{A}_i\) , \(\mathbf{E} = \varinjlim \mathbf{E}_i\) ； \(\mathbf{E}\) 具有 A-模结构，即诸 \(\mathbf{A}_i\) -模 \(\mathbf{E}_i\) 的结构的正向极限（II，§ 6，第 2 号）；此外，E 上的合成律——把诸 \(E_{i}\) 视为乘法下的广群（I，§ 10，第 3 号）而取其正向极限——连同 E 上的 A-模结构，在 E 上定义了一个 A-代数结构； \((\mathbf{E}_{i}, f_{jt})\) 称为 \(A_{i}\) -代数的正向系统，而 A-代数 E 称为其正向极限。若 \(f_{i}: E_{i} \to E\) , \(\phi_{i}: A_{i} \to A\) 为典范映射，则 \(f_{i}\) 是代数的 \(A_{i}\) -同态（相对于 \(\phi_{i}\) ）。若诸 \(E_{i}\) 结合（相应地，交换），则 E 亦然；若每个 \(E_{i}\) 具有单位元 \(e_{i}\) 且 \(f_{jt}(e_{i}) = e_{j}\) （对 \(i \leqslant j\) ），则 E 具有单位元 e，使得 \(f_{i}(e_{i}) = e\) 对所有 \(i \in I\) 成立。

设 \((\mathbf{E}_i', f_{ji}')\) 是 \(\mathbf{A}_i\) -代数的另一个正向系统，且对所有 \(i\) ，设 \(u_i: \mathbf{E}_i \to \mathbf{E}_i'\) 为 \(\mathbf{A}_i\) -代数同态，这些映射构成一个正向系统；则 \(u = \lim u_i\) 是一个 \(\mathbf{A}\) -代数同态。

现在假设所有的环 \(\mathbf{A}_i\) 都等于同一个环 \(\mathbf{A}\) ，且诸 \(\phi_{ji}\) 等于 \(\mathrm{Id}_{\mathbf{A}}\) ，于是 \(\mathbf{E} = \varinjlim \mathbf{E}_i\) 是一个 A-代数。设 \(\mathbf{F}\) 是一个 A-代数，且对所有 \(i\) ，设 \(u_i: \mathbf{E}_i \to \mathbf{F}\) 是一个 A-代数同态，使得 \((u_i)\) 是一个映射的正向系统；则 \(u = \varinjlim u_i\) 是代数 \(\mathbf{E}\) 到代数 \(\mathbf{F}\) 中的同态。反之，对每个 A-代数同态 \(v: \mathbf{E} \to \mathbf{F}\) ，族 \(v_i = v \circ f_i\) 是一个 A-代数同态的正向系统，使得 \(v = \varinjlim v_i\) 。此外，由于记 \(\bar{f}_{ij} = \operatorname{Hom}(f_{ij}, 1_F)\) , \((\operatorname{Hom}_{\mathbf{A-alg}}(\mathbf{E}_i, \mathbf{F}), \bar{f}_{ij})\) 显然是一个集合的逆向系统，可见上述评注也可以表述为：典范映射 \(v \mapsto (v \circ f_i)\) 是一个双射

\[
d _ {\mathrm{F}}: \operatorname{Hom} _ {\mathrm{A-alg.}} (\varinjlim \mathrm{E} _ {i}, \mathrm{F}) \rightarrow \varprojlim \operatorname{Hom} _ {\mathrm{A-alg.}} (\mathrm{E} _ {i}, \mathrm{F}).
\]

进一步地，对于每个A-代数同态 \(w: F \rightarrow F'\) ，

\[
\bar {w} _ {i} = \operatorname{Hom} (1 _ {\mathrm{E} _ {i}}, w): \operatorname{Hom} _ {\mathrm{A-alg.}} (\mathrm{E} _ {i}, \mathrm{F}) \rightarrow \operatorname{Hom} _ {\mathrm{A-alg.}} (\mathrm{E} _ {i}, \mathrm{F} ^ {\prime})
\]

构成映射的逆向系统，且图表

\[
\begin{array}{l} \operatorname{Hom} _ {\mathrm{A-alg.}} (\varinjlim \mathrm{E} _ {i}, \mathrm{F}) \xrightarrow {d _ {\mathrm{F}}} \varprojlim \operatorname{Hom} _ {\mathrm{A-alg.}} (\mathrm{E} _ {i}, \mathrm{F}) \\ \operatorname{Hom} (1 _ {\mathrm{E}}, w) \Biggl \downarrow \qquad \qquad \qquad \qquad \qquad \qquad \qquad \qquad \qquad \qquad \qquad \qquad \qquad \qquad \qquad \qquad \qquad \qquad \qquad \qquad \varprojlim \varpi_ {i} \\ \operatorname{Hom} _ {\mathrm{A-alg.}} (\varinjlim \mathrm{E} _ {i}, \mathrm{F} ^ {\prime}) \xrightarrow {d _ {\mathrm{F} ^ {\prime}}} \varprojlim \operatorname{Hom} _ {\mathrm{A-alg.}} (\mathrm{E} _ {i}, \mathrm{F} ^ {\prime}) \end{array}
\]

是交换的。

